\chapter{Pneumonectomy}

\section*{Introduction \& Clinical Indications}
Pneumonectomy is a radical surgical procedure that entails the complete excision of an entire lung, along with its associated hilar structures, including the main pulmonary artery, pulmonary veins, and main bronchus. This extensive resection is primarily indicated for malignant conditions such as centrally located non-small cell lung cancer (NSCLC) or malignant pleural mesothelioma, where the tumor extensively involves the hilum or main bronchus, precluding a more limited lung-sparing resection. Less commonly, pneumonectomy may be performed for severe, diffuse, and unilateral benign lung diseases like extensive bronchiectasis, chronic lung abscesses, or intractable fungal infections that are refractory to medical management and cause significant morbidity. Due to the profound physiological impact of losing an entire lung, patient selection is exceptionally rigorous, focusing on adequate cardiopulmonary reserve and the potential for curative intent, making it one of the most challenging operations in thoracic surgery.

\begin{itemize}
  \item Total lung resection
  \item Complete lung removal
  \item Standard pneumonectomy
  \item Extrapleural pneumonectomy (EPP)
  \item Intrapericardial pneumonectomy
  \item Radical pneumonectomy
\end{itemize}

The fundamental principle of pneumonectomy is the complete extirpation of diseased lung tissue, most often to achieve oncological clearance with negative surgical margins (R0 resection) for malignant tumors. This involves meticulous dissection and individual division of the hilar structures: the main pulmonary artery, the superior and inferior pulmonary veins, and the main bronchus, ensuring no residual tumor cells are left behind. The goal is to remove the entire lung en bloc, minimizing tumor spillage and ensuring radical resection, particularly for centrally located tumors or those involving the main bronchus or hilar vessels. 

For benign conditions, the rationale is to remove a source of intractable infection, hemorrhage, or severe functional impairment that cannot be managed by less extensive means. A critical technical goal is the secure closure and reinforcement of the bronchial stump to prevent a bronchopleural fistula (BPF), a potentially life-threatening complication. Following lung removal, the empty hemithorax undergoes a series of physiological adaptations: the mediastinum shifts towards the operative side, and the diaphragm elevates, reducing the volume of the pleural space. This space typically fills with serous fluid, which eventually organizes into a fibrous peel, effectively obliterating the cavity. The physiological consequence of pneumonectomy is a significant reduction in pulmonary function, as the remaining lung must compensate for the entire body's gas exchange. This necessitates careful preoperative assessment of predicted postoperative lung function and cardiac reserve to ensure the patient can tolerate the procedure and maintain an acceptable quality of life.

Early attempts at lung resection were fraught with prohibitively high mortality rates due to uncontrolled hemorrhage, infection, and respiratory compromise. The first successful one-stage pneumonectomy was performed by Dr. Evarts A. Graham in 1933 at Barnes Hospital in St. Louis, Missouri, for a 48-year-old physician with squamous cell carcinoma of the lung. This landmark operation, which involved the use of a cautery snare and mass ligation of the hilum, revolutionized thoracic surgery and demonstrated the feasibility of total lung removal. Prior to Graham's success, surgeons like Ferdinand Sauerbruch in Germany had performed partial resections, but a complete pneumonectomy was considered too dangerous.

The technique evolved significantly with advancements in anesthesia, particularly the development of endotracheal intubation and one-lung ventilation, which allowed for better control of the airway and lung collapse, improving surgical exposure and safety. The introduction of individual ligation and division of hilar structures by surgeons like Richard Overholt in the 1940s further improved safety and reduced complications compared to mass ligation. The advent of surgical staplers in the latter half of the 20th century provided more secure and reproducible methods for bronchial and vascular stump closure, significantly reducing the incidence of bronchopleural fistula and further refining the procedure, contributing to improved outcomes and making pneumonectomy a more standardized, albeit still high-risk, operation.

Pneumonectomy is a definitive surgical intervention ordered for specific, often severe, clinical indications, primarily for malignancy:

\begin{itemize}
  \item Centrally located non-small cell lung cancer (NSCLC) that involves the main bronchus, hilum, or pulmonary artery/veins, making a lobectomy or sleeve resection technically impossible or oncologically unsound to achieve clear margins.
  \item Extensive endobronchial tumors that obstruct the main bronchus and cannot be adequately resected by lesser means, leading to distal lung collapse or recurrent infections.
  \item Malignant pleural mesothelioma, often requiring an extrapleural pneumonectomy (EPP) as part of a multimodality treatment strategy, involving resection of the lung, pleura, pericardium, and diaphragm.
  \item Diffuse, unilateral, and severe benign lung diseases such as extensive bronchiectasis, chronic lung abscesses, or fungal infections (e.g., aspergilloma) that are refractory to medical therapy and cause significant morbidity like intractable hemoptysis or recurrent sepsis.
  \item Severe lung trauma with irreparable damage to the main bronchus or hilar vessels, leading to massive hemorrhage, complete lung destruction, or persistent air leak.
  \item Rarely, for extensive radiation necrosis, severe iatrogenic injuries to the lung, or congenital malformations causing intractable symptoms.
\end{itemize}

Pneumonectomy primarily serves as a definitive, potentially curative treatment for selected patients with centrally located lung malignancies. In this context, it is considered a primary surgical intervention aimed at achieving complete tumor eradication with negative surgical margins (R0 resection). For patients with malignant pleural mesothelioma, an extrapleural pneumonectomy (EPP) is often a primary component of a multimodality treatment regimen, combining aggressive surgery with chemotherapy and radiation therapy to improve survival. The decision to proceed with pneumonectomy in these cases is made after thorough multidisciplinary tumor board discussion, considering the patient's overall health, tumor stage, and potential for cure.

In certain benign conditions, such as severe, unilateral bronchiectasis, chronic lung abscesses, or intractable fungal infections, pneumonectomy is typically a secondary or salvage procedure. It is reserved for cases where medical management has failed, and less extensive surgical options (e.g., lobectomy, segmentectomy) are deemed insufficient to control the disease or are technically impossible due to the diffuse nature of the pathology. It is rarely a secondary procedure for failed primary lung cancer resections, as recurrence often indicates systemic disease that is better managed with systemic therapies. Given its significant morbidity and impact on quality of life, the decision to proceed with pneumonectomy is always made after thorough evaluation and consideration of all alternative treatments.

\section*{Relevant Surgical Anatomy}
A thorough understanding of thoracic anatomy is paramount for safe and effective pneumonectomy. The lungs are housed within the pleural cavities, separated by the mediastinum, which contains the heart, great vessels, trachea, esophagus, and major nerves. Each lung is supplied by a main bronchus, a main pulmonary artery, and two pulmonary veins (superior and inferior). The right main bronchus is shorter and wider than the left, branching into three lobar bronchi, while the left main bronchus is longer and narrower, branching into two lobar bronchi. 

The pulmonary arteries, arising from the right ventricle, carry deoxygenated blood to the lungs, typically lying superior and posterior to the main bronchus on the right, and superior and anterior on the left. The pulmonary veins, carrying oxygenated blood to the left atrium, are generally anterior and inferior to the bronchus and artery. 

The hilum of each lung, located on its mediastinal surface, is the entry and exit point for these structures, along with bronchial arteries (arising from the aorta or intercostal arteries), lymphatic vessels, and nerves. Key nerve relations include the phrenic nerve, which innervates the diaphragm and runs anterior to the lung hilum, and the vagus nerve, which passes posterior to the hilum and gives off the left recurrent laryngeal nerve looping under the aortic arch. Lymphatic drainage follows the bronchial and pulmonary arteries and veins to hilar (N1) and then mediastinal (N2) lymph nodes, which are critical for cancer staging. Surgical landmarks include the azygos vein arching over the right main bronchus and the aortic arch superior to the left main bronchus, both crucial for orientation during hilar dissection. The pericardium, a fibrous sac enclosing the heart, is also a critical structure, especially in cases requiring intrapericardial pneumonectomy or extrapleural pneumonectomy for mesothelioma.

\section*{Preoperative Workup \& Preparation}
Extensive preoperative evaluation is crucial for patients undergoing pneumonectomy due to the significant physiological impact of the procedure and its high morbidity and mortality.

\begin{itemize}
  \item \textbf{Comprehensive Medical History and Physical Examination:} To identify all comorbidities, assess overall health status, and evaluate functional capacity.
  
  \item \textbf{Pulmonary Function Testing (PFTs):} Including spirometry (FEV1, FVC), lung volumes, and diffusion capacity (DLCO). Crucially, predicted postoperative (PPO) FEV1 and DLCO are calculated, often requiring a quantitative ventilation-perfusion (V/Q) scan or computed tomography (CT) volumetry to accurately estimate the contribution of the lung to be resected. A PPO FEV1 and DLCO of greater than 40\% are generally considered acceptable for pneumonectomy, though some centers may accept lower values with careful patient selection.
  
  \item \textbf{Cardiac Evaluation:} Electrocardiogram (ECG), echocardiogram, and often a cardiac stress test (e.g., dobutamine stress echocardiogram or myocardial perfusion scan) to assess cardiac reserve, especially in patients with a history of coronary artery disease, arrhythmias, or significant risk factors.
  
  \item \textbf{Imaging Studies:} High-resolution chest CT with intravenous contrast for detailed tumor staging, assessment of hilar involvement, and evaluation of mediastinal structures. Positron Emission Tomography (PET-CT) for metabolic activity and detection of distant metastases. Brain MRI and bone scan may be performed for complete staging in oncology patients.
  
  \item \textbf{Mediastinal Staging:} Endobronchial ultrasound (EBUS) with fine-needle aspiration (FNA) or mediastinoscopy to accurately assess lymph node involvement (N2 disease), which is critical for determining resectability and prognosis.
  
  \item \textbf{Nutritional Assessment:} Optimization of nutritional status, as malnutrition can impair wound healing and immune function, increasing postoperative complications.
  
  \item \textbf{Smoking Cessation:} Patients are strongly advised to stop smoking at least 4-6 weeks prior to surgery to improve pulmonary function, reduce airway secretions, and decrease postoperative pulmonary complications.
  
  \item \textbf{Preoperative Education and Consent:} Thorough discussion of risks, benefits, alternatives, and expected postoperative course, including the significant impact on respiratory function.
  
  \item \textbf{Medication Review:} Adjustment or discontinuation of anticoagulants and antiplatelet agents according to institutional protocols. Management of chronic conditions like diabetes and hypertension.
  
  \item \textbf{Prophylactic Antibiotics:} Administered intravenously within 60 minutes prior to skin incision to reduce the risk of surgical site infection.
  
  \item \textbf{NPO Status:} Patients are instructed to be NPO (nil per os) for a specified period (typically 6-8 hours for solids, 2 hours for clear liquids) before surgery.
\end{itemize}

\textbf{Absolute Contraindications:}

\begin{itemize}
  \item Inadequate predicted postoperative pulmonary reserve (e.g., PPO FEV1 or DLCO < 30-40\%, depending on institutional guidelines), indicating the patient would not tolerate the loss of an entire lung and would likely experience severe respiratory failure.
  
  \item Severe, uncorrectable cardiac disease (e.g., severe pulmonary hypertension, unstable angina, recent myocardial infarction within 3-6 months, severe valvular disease, or uncontrolled arrhythmias) that significantly increases operative risk to an unacceptable level.
  
  \item Unresectable metastatic disease outside the ipsilateral hemithorax (M1b or M1c disease), indicating that surgery would not be curative and would only add morbidity without survival benefit.
  
  \item Extensive local invasion of vital mediastinal structures (e.g., aorta, superior vena cava, esophagus, trachea, vertebral bodies) that precludes complete tumor resection with negative margins (R0 resection).
  
  \item Active systemic infection or sepsis that has not been adequately treated.
\end{itemize}

\textbf{Relative Contraindications \& Precautions:}

\begin{itemize}
  \item Advanced age (typically > 75-80 years) and significant comorbidities (e.g., severe chronic obstructive pulmonary disease, renal failure requiring dialysis, liver cirrhosis with portal hypertension) increase the risk of complications and may necessitate a more conservative approach or careful risk-benefit analysis.
  
  \item Right-sided pneumonectomy is generally associated with a higher risk of complications, particularly post-pneumonectomy pulmonary edema (PPPE) and cardiac herniation, compared to left-sided pneumonectomy, due to the larger lung volume removed and the closer proximity of the heart and great vessels.
  
  \item Neoadjuvant therapy (chemotherapy or radiation) can increase the risk of bronchial stump complications (e.g., fistula) and may require careful surgical planning, meticulous technique, and robust bronchial stump reinforcement.
  
  \item Impaired functional status or poor performance status (e.g., ECOG > 2), indicating limited ability to recover from a major surgical insult.
  
  \item Significant weight loss or malnutrition, which can impair healing and increase susceptibility to infection.
  
  \item Previous thoracic surgery on the ipsilateral side, which may result in dense adhesions and increased technical difficulty.
\end{itemize}

\section*{Surgical Technique \& Procedure}
The procedure is performed under general anesthesia with selective one-lung ventilation, typically using a double-lumen endotracheal tube or a bronchial blocker, to collapse the operative lung and facilitate surgical access. The patient is positioned in a lateral decubitus position, with the operative side uppermost, ensuring adequate padding of pressure points.

\begin{itemize}
  \item \textbf{Incision:} The most common approach is a posterolateral thoracotomy, which provides excellent exposure to the hilum and mediastinum. The incision typically extends from the mid-scapular line to the anterior axillary line, usually through the fifth intercostal space. Muscle-sparing or anterolateral thoracotomy approaches may also be used depending on surgeon preference, patient anatomy, and the specific pathology.
  
  \item \textbf{Pleural Entry and Exploration:} The chest cavity is entered, and the lung is carefully inspected to confirm resectability, assess tumor extent, and rule out unforeseen metastatic disease within the pleural space or mediastinum. Adhesions are lysed to mobilize the lung.
  
  \item \textbf{Hilar Dissection and Vascular Ligation:} The pulmonary ligament is divided, and the pleura overlying the hilum is incised. The main pulmonary artery, superior pulmonary vein, and inferior pulmonary vein are individually identified, meticulously dissected free from surrounding tissues, ligated with heavy sutures (e.g., 2-0 silk) or divided using vascular staplers (e.g., linear cutting staplers with appropriate cartridge length), and then transected. The order of vascular division may vary, but typically the veins are addressed first to prevent engorgement of the lung.
  
  \item \textbf{Bronchial Dissection and Division:} The main bronchus is then carefully dissected, ensuring adequate length for a tension-free, healthy stump. Care is taken to avoid devascularization of the bronchial stump. The bronchus is typically divided using a surgical stapler (e.g., linear cutting stapler with a specific bronchial cartridge) to ensure a secure, airtight closure. The staple line is often oversewn with absorbable sutures for added security.
  
  \item \textbf{Bronchial Stump Reinforcement:} To significantly reduce the risk of bronchopleural fistula, the bronchial stump is almost always reinforced with a pedicled flap of well-vascularized tissue. Common choices include an intercostal muscle flap, pleura, pericardial fat pad, or a portion of the diaphragm, which are sutured over the staple line.
  
  \item \textbf{Lung Removal and Hemostasis:} The entire lung is removed from the chest cavity. Meticulous hemostasis is achieved throughout the operative field, ensuring no active bleeding from the chest wall, mediastinum, or hilar stumps.
  
  \item \textbf{Closure:} The chest wall is closed in layers, typically involving reapproximation of the ribs with pericostal sutures, closure of the muscle layers, subcutaneous tissue, and skin. Unlike most other lung resections, chest tubes are typically not placed after a standard pneumonectomy to allow the pleural space to fill with serous fluid and promote mediastinal shift towards the operative side, which helps stabilize the remaining structures and prevents overexpansion of the contralateral lung.
\end{itemize}

\section*{Complications \& Risks}
Pneumonectomy carries the highest operative mortality and morbidity among all lung resections due to its extensive nature and profound physiological impact.

\begin{itemize}
  \item \textbf{General Surgical Complications:}
  
  \begin{itemize}
    \item \textbf{Bleeding:} Intraoperative or postoperative hemorrhage, potentially requiring blood transfusions or reoperation. This can arise from hilar vessels, intercostal arteries, or chest wall.
    
    \item \textbf{Infection:} Wound infection, pneumonia in the remaining lung, empyema (infection of the pleural space), or mediastinitis.
    
    \item \textbf{Anesthesia Complications:} Reactions to medications, aspiration pneumonitis, stroke, myocardial infarction, deep vein thrombosis (DVT), pulmonary embolism (PE), or death.
    
    \item \textbf{Pain:} Significant postoperative pain requiring aggressive multimodal analgesia.
  \end{itemize}
  
  \item \textbf{Procedure-Specific Risks:}
  
  \begin{itemize}
    \item \textbf{Bronchopleural Fistula (BPF):} Leakage from the bronchial stump, a severe and life-threatening complication that can lead to empyema, sepsis, and aspiration into the remaining lung, often requiring reoperation and prolonged hospitalization.
    
    \item \textbf{Post-Pneumonectomy Pulmonary Edema (PPPE):} Non-cardiogenic pulmonary edema in the remaining lung, a life-threatening complication with high mortality, particularly after right pneumonectomy, often due to fluid overload or inflammatory mediators.
    
    \item \textbf{Cardiac Herniation:} Displacement of the heart through a defect in the pericardium (iatrogenic or congenital), leading to kinking of great vessels, hemodynamic instability, and potentially cardiac arrest.
    
    \item \textbf{Atrial Fibrillation:} The most common cardiac arrhythmia after thoracic surgery, often requiring medical management with beta-blockers or calcium channel blockers.
    
    \item \textbf{Respiratory Failure:} Due to inadequate pulmonary reserve, pneumonia, acute lung injury in the remaining lung, or exacerbation of underlying lung disease, potentially requiring prolonged mechanical ventilation and tracheostomy.
    
    \item \textbf{Post-Pneumonectomy Syndrome:} A rare, late complication where the mediastinum shifts excessively, causing torsion and compression of the remaining main bronchus or great vessels, leading to respiratory distress or vascular compromise.
    
    \item \textbf{Stump Dehiscence:} Breakdown of the bronchial or vascular stump closure, leading to BPF or hemorrhage.
    
    \item \textbf{Chylothorax:} Leakage of lymphatic fluid into the chest cavity if the thoracic duct or its major branches are injured during mediastinal dissection, requiring drainage and potentially reoperation.
    
    \item \textbf{Nerve Injury:} Damage to the recurrent laryngeal nerve (leading to vocal cord paralysis and hoarseness) or phrenic nerve (leading to ipsilateral diaphragmatic paralysis and reduced respiratory efficiency).
    
    \item \textbf{Acute Kidney Injury:} Can occur due to hypoperfusion during surgery, nephrotoxic medications, or systemic inflammatory response.
  \end{itemize}
\end{itemize}

\section*{Postoperative Care \& Recovery}
Immediately after pneumonectomy, the patient is transferred to a specialized intensive care unit (ICU) or a high-dependency unit for close monitoring. Pain management is critical and often involves epidural analgesia, patient-controlled analgesia (PCA) with intravenous opioids, or regional nerve blocks. Respiratory status, oxygen saturation, heart rate, blood pressure, central venous pressure, and fluid balance are meticulously monitored. Unlike other lung resections, chest drains are typically avoided in standard pneumonectomy to allow the empty pleural space to fill with fluid, which helps stabilize the mediastinum and prevent excessive shift.

Over the first 24-48 hours, the focus is on maintaining adequate oxygenation, aggressive pain control, and preventing complications. Early mobilization, often within 24 hours, is strongly encouraged to reduce the risk of deep vein thrombosis and pulmonary embolism. Respiratory physiotherapy, including incentive spirometry and deep breathing exercises, is initiated to promote lung expansion of the remaining lung and clear secretions. Fluid management is particularly delicate to avoid both dehydration and fluid overload, which can precipitate post-pneumonectomy pulmonary edema. Patients typically remain in the hospital for 7-14 days, depending on their recovery trajectory and the absence of complications. The first 1-2 weeks involve gradual increases in activity, a progressive diet, and careful monitoring for signs of infection or BPF. Long-term recovery involves a gradual return to full activity over several months, often with pulmonary rehabilitation, and adaptation to the reduced lung capacity, which may manifest as persistent dyspnea on exertion.

During the pneumonectomy, the surgeon documents key intraoperative findings, including the precise location and extent of the primary tumor, its relationship to hilar structures, the presence of any pleural effusions or adhesions, and the status of regional lymph nodes. Any suspicious areas are biopsied for frozen section analysis. The resected lung, along with its hilar structures and any associated lymph nodes, is sent for detailed pathological examination.

The pathology report is crucial, especially for malignancy. It provides definitive information on the tumor type (e.g., adenocarcinoma, squamous cell carcinoma, mesothelioma), histological grade, size, and extent of invasion into surrounding tissues (e.g., visceral pleura, chest wall, pericardium). Most importantly, it confirms the status of the surgical margins (negative margins, or R0 resection, indicate complete removal of visible tumor; positive margins, or R1/R2 resection, indicate residual tumor). The report also details the number and location of involved lymph nodes (N-stage), which is vital for accurate cancer staging. For benign conditions, the pathology report confirms the diagnosis (e.g., extensive bronchiectasis, chronic fungal infection) and the complete removal of the diseased tissue. These findings guide decisions regarding adjuvant therapies and inform the long-term prognosis.

A normal, successful postoperative course following pneumonectomy is characterized by a stable and uncomplicated recovery. Patients typically experience significant but manageable pain, which gradually subsides over days to weeks with appropriate analgesia. Respiratory function, while reduced, should be adequate, allowing for successful weaning from ventilatory support and maintenance of acceptable oxygen saturation on room air or with minimal supplemental oxygen. The mediastinum should remain stable, with a gradual shift towards the operative side as the pleural space fills with fluid.

Clinically, the patient should experience resolution of the symptoms that necessitated the surgery, such as intractable hemoptysis, recurrent infections, or severe dyspnea caused by the diseased lung. Wound healing should progress without infection or dehiscence. Over the long term, a normal course involves successful adaptation to the loss of one lung, with the patient achieving an acceptable quality of life and functional status. While some degree of dyspnea on exertion is common, patients should be able to perform activities of daily living and gradually increase their exercise tolerance, often with the aid of pulmonary rehabilitation. The absence of major complications like bronchopleural fistula, empyema, or post-pneumonectomy pulmonary edema is the hallmark of a successful outcome.

Post-pneumonectomy follow-up is critical for monitoring recovery, detecting complications, and surveilling for disease recurrence, especially in oncology patients.

\begin{itemize}
  \item \textbf{Hospital Discharge Planning:} Detailed instructions on wound care, pain management, activity restrictions (e.g., avoiding heavy lifting for 6-8 weeks), and warning signs (e.g., fever, increasing shortness of breath, new cough, wound drainage, chest pain) requiring immediate medical attention.
  
  \item \textbf{Postoperative Clinic Visits:} Typically scheduled at 1-2 weeks for wound check and suture/staple removal, and then at regular intervals (e.g., 1, 3, 6, 12 months, then annually) to assess recovery, functional status, and address any concerns.
  
  \item \textbf{Imaging Surveillance:} For oncology patients, regular chest CT scans (e.g., every 6-12 months for several years) are performed to monitor for local recurrence in the mediastinum or distant metastases.
  
  \item \textbf{Pulmonary Function Tests (PFTs):} May be repeated periodically to assess the long-term adaptation and function of the remaining lung and guide rehabilitation efforts.
  
  \item \textbf{Cardiac Evaluation:} Ongoing monitoring for cardiac arrhythmias or other cardiovascular issues, especially in patients with pre-existing cardiac disease.
  
  \item \textbf{Pulmonary Rehabilitation:} Often recommended to help patients improve their exercise tolerance, breathing techniques, and overall functional capacity through structured exercise and education programs.
  
  \item \textbf{Oncology Follow-up:} For cancer patients, regular visits with a medical or radiation oncologist to discuss adjuvant therapy, monitor tumor markers, and manage long-term cancer care.
\end{itemize}

Patients are advised to report any new or worsening symptoms, such as persistent cough, increasing shortness of breath, chest pain, unexplained weight loss, or changes in voice, to their healthcare provider promptly.

\section*{Outcomes \& Prognosis}
Pneumonectomy is a relatively uncommon surgical procedure compared to lobectomy, reflecting its radical nature and the stringent selection criteria. The incidence of pneumonectomy has been declining over recent decades due to advances in diagnostic imaging allowing for earlier detection of smaller, more peripheral tumors, and the development of lung-sparing techniques like sleeve resections and segmentectomies. However, it remains an essential procedure for specific indications where less extensive resections are not feasible. The vast majority of pneumonectomies (over 80-90\%) are performed for malignant conditions, primarily non-small cell lung cancer (NSCLC), with a smaller proportion for malignant pleural mesothelioma or severe benign diseases.

The procedure is more frequently performed in older adults, typically between 60 and 80 years of age, reflecting the peak incidence of lung cancer. There is no significant sex predilection for pneumonectomy itself, but the underlying disease epidemiology (e.g., historically higher lung cancer rates in men, though rates are converging) may influence patient demographics. Operative mortality rates for pneumonectomy for NSCLC range from 3\% to 10\% in experienced centers, with right pneumonectomy generally carrying a higher risk (5-15\%) than left pneumonectomy (3-8\%) due to the larger lung volume removed and the closer proximity of the heart. Morbidity rates are substantial, often exceeding 30-40\%, with complications such as atrial fibrillation, bronchopleural fistula, and post-pneumonectomy pulmonary edema being significant contributors to prolonged hospital stays and increased healthcare costs.

The outcomes and prognosis following pneumonectomy are highly dependent on the underlying indication, the patient's preoperative health status, the presence of complications, and the pathological findings. For non-small cell lung cancer, the 5-year survival rates after pneumonectomy vary significantly with the pathological stage of the disease. For Stage I NSCLC, 5-year survival can be around 40-50\%, decreasing to 20-30\% for Stage II, and even lower for Stage III disease, particularly with N2 nodal involvement. Achieving clear surgical margins (R0 resection) is the most critical prognostic factor for long-term survival in cancer patients, often necessitating a multidisciplinary approach involving adjuvant chemotherapy or radiation therapy.

Functional outcomes are also a major consideration. Patients typically experience a significant reduction in exercise tolerance and may have persistent dyspnea, especially with exertion, due to the loss of an entire lung. Quality of life can be impacted, though many patients adapt well over time, particularly with dedicated pulmonary rehabilitation programs. Recurrence rates for lung cancer after pneumonectomy are influenced by tumor stage, nodal status, and the completeness of resection, with local or distant recurrence occurring in a substantial proportion of patients, highlighting the importance of long-term surveillance. For malignant pleural mesothelioma treated with extrapleural pneumonectomy, the prognosis remains guarded, with median survival typically ranging from 12 to 24 months, despite aggressive multimodality treatment. For benign conditions, the prognosis is generally excellent once the acute postoperative period is navigated, with complete resolution of the underlying disease and a good long-term quality of life.

\section*{References}
\begin{enumerate}
  \item MedlinePlus. Pneumonectomy. U.S. National Library of Medicine; 2024. Available from: \url{https://medlineplus.gov/ency/article/002953.htm}
  
  \item Sabiston Textbook of Surgery: The Biological Basis of Modern Surgical Practice. 22nd ed. Elsevier; 2022. Chapter 56: Lung, Chest Wall, Pleura, and Mediastinum.
  
  \item National Comprehensive Cancer Network (NCCN) Clinical Practice Guidelines in Oncology (NCCN Guidelines\textregistered) for Non-Small Cell Lung Cancer. Version 5.2024.
  
  \item Shields' General Thoracic Surgery. 8th ed. Wolters Kluwer; 2019. Chapter 46: Pneumonectomy.
\end{enumerate}

\clearpage
